% preamble.tex — packages, figure toolchain, and the citation macros.
% Loaded by main.tex before notation.tex. No chapter content lives here.
%
% CITATION SCHEME (pinned by bp-077 §6, Q3 — one mechanism for the whole book):
%   \artifact{<id>}        design claim  -> a ratified design-note / finding id
%   \coderef{<path>}{<ref>} implementation claim -> code by path @ git-ref
% Snippets are COPIES annotated `source: path@ref`. Draft notes are barred as
% citations (dn-agent-workflow §3); draft theses are forward-\ref'd, never cited.

\usepackage[utf8]{inputenc}
\usepackage[T1]{fontenc}
\usepackage{lmodern}
\usepackage[a4paper,margin=1in]{geometry}
\usepackage{microtype}
\usepackage{amsmath,amssymb,amsthm}
\usepackage{xcolor}
\usepackage{enumitem}
\usepackage{booktabs}   % added by Ch.2 (bp-104): rules for the tier/zone tables
\usepackage{tikz}
\usetikzlibrary{positioning,arrows.meta,calc,fit,backgrounds,shapes.geometric}
% hyperref last (except for the macros it must see); cleveref would go after it.
\usepackage[hidelinks]{hyperref}

% --- The edition marker: the HEAD sha this edition is keyed to (see SYNC.md). ---
% Every \coderef in the book pins to this ref unless a snippet states otherwise.
% bp-104 re-keyed bdcd9bc -> 009b726 after re-verifying EVERY citation and snippet
% against that ref, file by file (260 commits had passed; spot-checking was refused).
\newcommand{\gitref}{009b726}

% --- Citation macros (bp-077 §6). ------------------------------------------------
% \artifact — a design-note or finding id, rendered as a small-caps bracketed token.
\newcommand{\artifact}[1]{\textnormal{\textsc{[#1]}}}
% \coderef — implementation reference: <path> @ <git-ref>, both monospaced.
% Escape underscores in the path argument at the call site (e.g. check\_imports.py).
\newcommand{\coderef}[2]{\mbox{\texttt{#1}\,@\,\texttt{#2}}}

% --- Named colours for the two ratified diagrams. --------------------------------
\definecolor{inkblue}{HTML}{1F3A5F}
\definecolor{gate}{HTML}{B23A48}
\definecolor{soft}{HTML}{6B7B8C}
\definecolor{panel}{HTML}{EEF2F6}

% --- Book-specific environments. -------------------------------------------------
% A "Principle" is a stated stance of the system, always followed by its citation.
\theoremstyle{definition}
\newtheorem{principle}{Principle}[chapter]
% A "design-evolution remark" carries superseded material as pedagogy, warrant-linked
% (book skill §chapter-map). First used in Ch.2; defined once, here.
\newtheorem*{devolution}{Design-evolution remark}
% Added by Ch.2 (bp-104). An "Invariant" is a property the system enforces, stated
% formally and always cited to the ratified source that decides it. A "Proposition"
% is a consequence the manual DERIVES from a cited invariant — its proof is given
% inline and is checkable by the reader; it asserts nothing the record does not.
\theoremstyle{definition}
\newtheorem{invariant}{Invariant}[chapter]
\theoremstyle{plain}
\newtheorem{proposition}{Proposition}[chapter]
% hyperref needs an \autoref name for every custom theorem environment, else it
% warns and prints nothing before the number.
\providecommand{\invariantautorefname}{Invariant}
\providecommand{\propositionautorefname}{Proposition}
\providecommand{\principleautorefname}{Principle}

% Draft-thesis forward-reference helper: names a thesis whose formal design note is
% still draft (barred from the book, dn-agent-workflow §3) and points at the chapter
% that will carry it once ratified. Keeps the "no draft citation" rule mechanical.
% Usage: \fwdthesis{authorship distance}{ch:architecture}
\newcommand{\fwdthesis}[2]{\emph{#1}~(design in progress; its formal treatment lands
  in \autoref{#2} on ratification)}
